\documentclass{article}
\usepackage[T1]{fontenc}
\usepackage{amsmath}
\begin{document}
\section*{Issue 92: finite distributor grids and uniform continuation}
\paragraph{Target and hypotheses.}
The source is Sections 2--3 of
\texttt{General\_Independent\_Interaction\_Criterion.tex}, in particular
the addressed regular-distributivity step, the constant-on-a-fibre step
in Lemma~\texttt{lem:depfunctor}, and the structural normalization step
of Proposition~\texttt{prop:fubini}.
The coefficient category is any \(C:\mathrm{Type}_u\) with morphisms in
universe \(v\). The theorem uses two bifunctors, an additive unit,
a natural invertible middle-four arrow, an additive merge of two zeros,
actual distributors and absorbers, and their binary and nullary
compatibility equations. Row and column naturality additionally uses
the binary distributor naturality and nullary absorber equations.
\texttt{Additive.ofBraided} constructs the additive input from a braided
monoidal structure; \texttt{Binary.ofMonoidal} gives the bifunctor of any
supplied monoidal structure. The distributors and their compatibility
remain explicit inputs. No coproduct, biproduct, or initiality is used.

\paragraph{Checked claims.}
\texttt{PreservesSplit.fold} inducts from the binary and nullary laws to
arbitrary right folds. \texttt{GridData.interchange} identifies the right-first
and left-first coefficient expansions through an independent recursive
transpose; \texttt{GridData.transpose\_natural},
\texttt{GridNatural.rows} and \texttt{GridNatural.cols} give commutation
on actual arrows. The transpose uses the canonical middle swap at each
two-row step; repeated equal object entries can still be distinct
occurrences. The lists may be empty and their lengths have no bounds.
The type \texttt{Arrows} records an actual arrow at each occurrence.
\texttt{GridData.constantOnFibre} uses an arbitrary map
\(\bigoplus_x V_x \to A\) from the entire sum. A common continuation
\(B\to D\) commutes with the merge after the specified inverse distributor.
The proof uses structural naturality and tensor bifunctoriality, including
the empty fibre. A different continuation at each input is not admitted
by this local interchange law.

\paragraph{Instance and scope.}
The probe category is the opposite of small types, presented by reverse
functions, with disjoint sum and cartesian product as the two object
operations and with all needed diagrams proved directly. The stated
probes use finite types including an empty type. An actual absent
coproduct injection shows why the proof cannot rely on coproducts.
An explicit reverse function into a tagged sum is used as an opaque
aggregate. A computed entry and repeated-coefficient transpose verify
the tags, and zero-size cases are separate.
The construction treats rectangular right-fold grids and a single normalized
outer fibre. It does not construct reindexing along arbitrary finite carrier
maps, partition/gluing coherence, augmented-family category and terminal
object, arbitrary dependent accumulation, addressed independent tensor
coherence, or the packaged augmented Fubini theorem. Those are open at this
generality. The existing \#241 matrix instance is unchanged.
\end{document}
